\section{Claim-by-claim coverage of the source page}\label{app:coverage}
The following concordance refers to revision 1368909113 of \cite{wiki}. It
tracks mathematical content rather than historical prose; repeated displays
of the same identity are grouped, but different hypotheses and convergence
interpretations are distinguished. An entry marked \emph{corrected} or
\emph{qualified} gives the valid statement and identifies why the
unqualified original cannot be asserted; it does not mean that a false
unrestricted assertion has been proved. Theorems appearing only as names
in links are not statements made by the page; the Trotter and higher-order
splitting results are included as additional consequences.

\begingroup\small
\setlength{\tabcolsep}{4pt}
\renewcommand{\arraystretch}{1.18}
\begin{longtable}{@{}>{\raggedright\arraybackslash}p{0.21\textwidth}>{\raggedright\arraybackslash}p{0.40\textwidth}>{\raggedright\arraybackslash}p{0.34\textwidth}@{}}
\caption{Claim-by-claim mathematical coverage.}\label{tab:coverage}\\
\toprule
Source location & Formula or assertion & Proof or qualification here\\
\midrule\endfirsthead
\multicolumn{3}{c}{\tablename\ \thetable\ (continued)}\\
\toprule
Source location & Formula or assertion & Proof or qualification here\\
\midrule\endhead
\bottomrule\endfoot
Introduction & Exponential and logarithm series; $e^Xe^Y=e^Z$ with $Z$ a formal series of iterated commutators; leading terms. & \Cref{lem:explog,thm:wordcut,thm:formalbch,thm:dynkin}; \eqref{eq:z1}--\eqref{eq:z3}.\\
Introduction & Commutator interpretation, antisymmetry, Jacobi identity. & \Cref{sec:scope}, \eqref{eq:derivation}, by associative expansion.\\
Introduction & Small elements give a convergent BCH series; the exponential gives coordinates near the identity; local multiplication is determined by the bracket. & \Cref{thm:convergence,thm:local}.\\
Introduction; explicit forms & Lie correspondence; integration of Lie algebra homomorphisms and representations. & \Cref{thm:localcorr,thm:integration}; simple connectedness stated as the global hypothesis.\\
Dynkin formula & Complete block sum, factorial and total-degree denominators, right-nested convention. & \Cref{thm:dynkin}, \eqref{eq:dynkin}, \eqref{eq:rightbracket}.\\
Dynkin formula & Vanishing of summands whose final block ends in repeated letters. & Last assertion of \Cref{thm:dynkin}, $[a,a]=0$.\\
Dynkin formula & Every displayed BCH term through degree six. & \eqref{eq:z1}--\eqref{eq:z6}, \Cref{prop:lowdegree}, certificate \Cref{tab:words6}.\\
Dynkin formula & All nonlinear terms lie in the ideal generated by $[X,Y]$. & \Cref{thm:formalbch}.\\
Dynkin formula & $Z(Y,X)=-Z(-X,-Y)$; homogeneous interchange signs. & \Cref{prop:symmetry}, \eqref{eq:swap}, \eqref{eq:parity}.\\
Integral form & Poincar\'e integral and the function $\psi(x)=x\log x/(x-1)$. & \Cref{thm:poincare}, \eqref{eq:poincare}, \eqref{eq:psi}.\\
Integral form; note & $\psi(e^y)$ and the Bernoulli coefficients with $B_1=+\tfrac12$. & \eqref{eq:psibeta}, \Cref{prop:bernoulli}.\\
Matrix illustration & Tangent Lie algebra, matrix commutator, matrix exponential. & Opening of \Cref{sec:liegroups}.\\
Matrix illustration & Full associative logarithmic block sum. & \eqref{eq:blocks}, \Cref{thm:wordcut}; absolute convergence in \Cref{thm:convergence}.\\
Matrix illustration & Associative terms $z_1,\ldots,z_4$ and their expression by brackets. & \eqref{eq:assoc2}--\eqref{eq:assoc4}, \Cref{prop:lowdegree}.\\
Matrix illustration; note & Conditions $\norm X+\norm Y<\log2$ and $\norm Z<\log2$; Hilbert--Schmidt norm. & \Cref{thm:convergence}; discussion after \eqref{eq:mercatortail}; submultiplicativity of the HS norm proved in \Cref{sec:convergence}.\\
Matrix illustration & Trace of the BCH logarithm; higher terms traceless. & \Cref{prop:trace}; \emph{qualified} for other branches by \eqref{eq:tracebranch}.\\
Convergence issues & The $\mathfrak{sl}_2(\C)$ matrices, their product, and absence of a traceless logarithm; divergence of BCH. & \Cref{thm:sl2}; all logarithms classified, even conditional convergence excluded.\\
Convergence issues & Nonisomorphic groups with the same Lie algebra. & The $\R$ versus $S^1$ example after \Cref{thm:local}.\\
Convergence issues & Bound $\norm X+\norm Y<\tfrac12\log2$ in any submultiplicative norm. & \Cref{thm:convergence}(iii), distinguished from the grouped bound.\\
Special cases & Commuting variables. & \Cref{prop:commuting}.\\
Special cases; Campbell application & Central commutator: truncation and disentangling without smallness. & \Cref{thm:central}, including the differential proof of global validity.\\
Special cases & Expansion linear in $Y$, hyperbolic cotangent form, Bernoulli coefficients of $\ad_X^nY$. & \Cref{cor:linearY}; all higher $Y$-orders in \eqref{eq:bidegree}, \eqref{eq:Hqrec}, \eqref{eq:H2kernel}.\\
Special cases & $[X,Y]=sY$: no terms with two or more $Y$'s; closed form; excluded poles; removable $s=0$. & \Cref{thm:eigen}, \Cref{ex:resonance}; Taylor convergence versus global resummation distinguished.\\
Special cases & Eigenvector braiding and adjoint dilation. & \eqref{eq:eigenbraid}, \Cref{thm:campbell}.\\
Existence results & Formal exp/log bijection; homogeneous grouping. & \Cref{lem:explog,thm:wordcut}.\\
Existence results & Coproduct on noncommuting series. & \eqref{eq:coproduct}, \eqref{eq:unshuffle}; \emph{qualified}: completed tensor product, with the counterexample in \Cref{app:pbw}.\\
Existence results & Exponential is group-like iff exponent is primitive; group-like elements form a group. & \Cref{lem:grouplike}, with constant term $1$ in the definition.\\
Existence results & Friedrichs' characterization of primitive elements; BCH follows. & \Cref{lem:dsw,thm:friedrichs,thm:formalbch}.\\
Existence results & Free enveloping algebra is the free associative algebra; natural Hopf structure; validity in any characteristic-zero Lie algebra. & \Cref{prop:enveloping}, \Cref{sec:universal}, \Cref{app:pbw}.\\
Existence results & Existence of a recursive Lie-only construction. & \Cref{thm:homrec}, independent of the Hopf proof.\\
Zassenhaus formula & Ordered product; displayed $C_2,C_3,C_4$ including the left-nested $C_4$. & \Cref{thm:zassenhaus}, \Cref{prop:zassdisplayed}.\\
Zassenhaus formula & Compact product $e^{-tX}e^{t(X+Y)}=\prod e^{t^nC_n}$. & Convention $C_1=Y$, after \Cref{thm:zassenhaus}.\\
Zassenhaus formula & All later exponents are homogeneous Lie polynomials; recursive determination; the ellipsis. & \Cref{thm:zassenhaus}, \Cref{prop:zassdiff}, \Cref{thm:tree}, \Cref{tab:zass56}; convergence in \Cref{thm:zconv}.\\
Zassenhaus formula & Suzuki--Trotter consequence. & \Cref{sec:splitting}: \Cref{thm:trotter,prop:symmetric,thm:suzuki} with complete error series.\\
Campbell identity & $\Ad_{e^X}=e^{\ad_X}$, all iterated commutators, and the defining differential equation. & \Cref{thm:campbell}, \eqref{eq:adbinomial}.\\
Campbell identity & General conjugation and braiding identities; central reduction. & \eqref{eq:conjexp}, \eqref{eq:braiding}, \eqref{eq:centralconj}.\\
Campbell application & Differentiation of $e^{sX}e^{sY}$ and its solution. & Proof of \Cref{thm:central}.\\
Infinitesimal case & Full series for $e^{-X}\dd e^X$. & \eqref{eq:leftdexp}, \eqref{eq:infinitesimal}.\\
Infinitesimal case & Structure-constant expansion; matrices $M$ and $W$; quotient notation. & \eqref{eq:components}--\eqref{eq:thetaW}; \emph{qualified}: $(I-e^{-M})M^{-1}$ means $\phi(M)$.\\
Infinitesimal case & Identification of $M$ as a vielbein. & \emph{Corrected} by \Cref{prop:coframe}: $W$ is the coframe and $M$ is always singular.\\
Infinitesimal case & Pullback metric $g_{ij}=W^a{}_iW^b{}_jB_{ab}$ and the Killing form as metric. & \Cref{thm:pullback}, \eqref{eq:killing}; \emph{qualified}: Jacobian factors and nondegeneracy required; all orders in \Cref{thm:metric}.\\
Quantum mechanics & $[Q,P]=\ii\hbar I$ and centrality. & \eqref{eq:QP}, \eqref{eq:CCR} on the stated invariant core.\\
Quantum mechanics & Exponentiated canonical relation; exact symmetric three-factor form. & \Cref{thm:weyl}; \Cref{prop:ccrcounter} shows a bare dense-domain commutator is insufficient.\\
Quantum mechanics & $[a^\dagger,a]=-I$; normal ordering of the displacement operator. & \eqref{eq:creation}, \Cref{thm:displacement}, with explicit domain and closure.\\
Quantum mechanics & Product of two displacements, its phase, Heisenberg nilpotence. & \Cref{thm:Dcomposition}, \eqref{eq:Dproduct}, \eqref{eq:heisenberglaw}.\\
Quantum mechanics & Expansions hidden in the displacement exponentials. & \eqref{eq:Dall}--\eqref{eq:laguerre}.\\
\end{longtable}
\endgroup

\subsection*{Summary of essential qualifications}
The tensor product in the formal existence proof must be completed, and
group-like elements have constant term one. The trace identity refers to the
BCH logarithm, not to an arbitrary branch. Smallness is a sufficient
convergence condition, not a necessary one, and resummed validity of a
closed form does not imply convergence of its Taylor series; at nonzero
resonances $2\pi\ii k$ the closed form of shape $X+cY$ fails although a
logarithm exists. In exponential coordinates $M=\ad_X$ is singular and is
not the coframe; $W=\phi(M)$ is the relevant matrix, and its inverse exists
exactly off the resonances. Neither an arbitrary pullback nor the Killing
form is automatically a metric. The alternate Zassenhaus product starting
at $n=1$ needs $C_1=Y$. Finally, an unbounded commutator identity requires
domain and exponentiation hypotheses; concrete realizations, not formal
substitution alone, justify the quantum formulas.

\subsection*{Beyond the finite displays}
Every BCH coefficient is given by \eqref{eq:wordcut} and
\eqref{eq:bidegreeproj}; every bidegree by \eqref{eq:bidegree}; every power
of $Y$ by \eqref{eq:Hqrec} and \eqref{eq:H2kernel}; every Zassenhaus
exponent by \eqref{eq:treeformula}; every Maurer--Cartan component by
\eqref{eq:MW}--\eqref{eq:thetaW}; every invariant-metric coefficient by
\eqref{eq:metriceven}; every splitting error by \eqref{eq:multibch}; and
every oscillator matrix element by \eqref{eq:Dmatrix}. These statements
specify coefficients at arbitrary order; the accompanying tables are finite
illustrations of unlimited formulas, not their claimed endpoint.
